\documentclass[11pt]{article}

% Compile with pdfLaTeX. Keep the qf_results directory beside this file.
\usepackage[T1]{fontenc}
\usepackage[utf8]{inputenc}
\usepackage{mathptmx}
\usepackage{courier}
\usepackage{textcomp}
\usepackage[margin=1in]{geometry}
\usepackage{graphicx}
\usepackage{booktabs}
\usepackage{array}
\usepackage{tabularx}
\usepackage{adjustbox}
\usepackage{caption}
\usepackage{amsmath}
\usepackage{amssymb}
\usepackage{enumitem}
\usepackage{microtype}
\usepackage[section]{placeins}
\usepackage{xurl}
\usepackage[hidelinks]{hyperref}

\captionsetup{font=small,labelfont=bf}
\setlength{\parindent}{0pt}
\setlength{\parskip}{0.5em}
\setlength{\emergencystretch}{2em}
\setlist[itemize]{leftmargin=*,itemsep=0.25em,topsep=0.25em}
\renewcommand{\arraystretch}{1.12}
\newcolumntype{Y}{>{\raggedright\arraybackslash}X}
\DeclareMathOperator{\sgn}{sign}

% Show a labelled placeholder if an external image is unavailable.
% Existing images are included automatically at their original paths.
\newcommand{\reportfigure}[1]{%
  \IfFileExists{#1}{%
    \includegraphics[width=\linewidth,height=0.72\textheight,keepaspectratio]{#1}%
  }{%
    \PackageWarning{reproduction}{Missing figure: #1}%
    \fbox{\parbox[c][1.15in][c]{0.94\linewidth}{%
      \centering\textbf{Figure file not found}\par
      \small\nolinkurl{#1}}}%
  }%
}

\title{Query-Free Adversarial Attacks on Stable Diffusion:\\ A Small-Scale Reproduction}
\author{S. M. Raihanul Alam\\ Department of Computer Science, University of Iowa}
\date{November 2026}

\begin{document}
\maketitle

\begin{abstract}
Zhuang et al.~\cite{zhuang2023} showed in his paper that appending 5 characters to a prompt can change the output of Stable Diffusion, using a search that only evaluates the CLIP text encoder and never calls the diffusion model. This report reimplements their untargeted and targeted objectives with \textbf{genetic}, \textbf{greedy}, and \textbf{gradient-based (PGD) search} algorithms and \textbf{includes a random-search baseline having the same evaluation budget as the genetic algorithm.}. On 2 prompts from the paper, all searches lowered the text-embedding objective well below a single random suffix for the same scene and objective, with cosine similarities of 0.52 to 0.77 against 0.76 to 0.86. The effect on images was much weaker. Mean CLIP score drops were at most 0.044, some were negative, and a separate object detector still found the target object in all eight targeted images. The genetic algorithm did not consistently beat the random search. An analysis of the embeddings help explain why the attack may fail to remove objects. The encoder output for the target word barely changes, over 99\% of the mask falls on padding positions, and every suffix pushes the embedding further from the version of the prompt without the target object. Since we tested only 2 prompts with one seed, these findings apply only to this run and cannot tell us whether the paper’s ranking of methods holds.
\end{abstract}

\section{Introduction}

Text-to-image models take their conditioning from a text encoder, so any weakness in that encoder may carry through to the generated image. Zhuang et al.~\cite{zhuang2023} use this to attack Stable Diffusion v1.4~\cite{rombach2022} with a short character suffix, chosen by querying only the CLIP text encoder~\cite{radford2021}. Because the encoder is public and cheap to run, the attacker needs no access to the image model during the search.

The paper makes 2 claims that this report examines. The first is that a 5 character suffix substantially changes the generated image. The second is that restricting the search to selected embedding dimensions can remove a chosen object while leaving the rest of the scene intact.

This work reproduces the method on a smaller scale and adds a random-search baseline with the same number of encoder evaluations as the genetic algorithm. This comparison helps show whether the genetic algorithm benefits from its search strategy or simply from trying many suffixes.

\section{Background}

\subsection{Threat model}

The attacker can evaluate the CLIP text encoder and can append up to 5 characters to a prompt. ``Query-free'' applies to the diffusion pipeline during the search only, so encoder evaluations are still required. The attack does not perturb image inputs, and the paper does not study safety-filter bypass.

% \subsection{Attack objectives}

% Let $\tau(x)$ be the text embedding of prompt $x$, and let $x'$ be the prompt with a suffix appended. The untargeted attack minimizes
% \begin{equation}
%   \mathcal{L}_{\mathrm{full}}(x,x') = \cos\bigl(\tau(x),\tau(x')\bigr).
%   \label{eq:full}
% \end{equation}
% The targeted attack applies the same loss to a subset of dimensions:
% \begin{equation}
%   \mathcal{L}_{\mathrm{masked}}(x,x') =
%   \cos\bigl(I\odot\tau(x),I\odot\tau(x')\bigr),
%   \label{eq:masked}
% \end{equation}
% where $I$ is a binary mask and $\odot$ denotes element-wise multiplication. To build the mask, $n$ sentences are encoded with and without the target phrase, giving difference vectors $d_i$. Dimension $j$ is selected when
% \begin{equation}
%   \left|\sum_{i=1}^{n}\sgn(d_{i,j})\right| > \varepsilon n.
%   \label{eq:mask}
% \end{equation}
% The paper uses $n=10$ and $\varepsilon=0.9$, so all ten signs must agree. Nothing in this loss explicitly preserves the rest of the scene.

\subsection{Published results}

Table~\ref{tab:published} summarizes the paper's results for 20 prompts
with 10 images each. Compared with no attack, genetic search reduces
the mean CLIP score from 0.277 to 0.203 in the untargeted setting
(a 26.7\% reduction) and from 0.229 to 0.186 in the targeted setting
(an 18.8\% reduction).

\begin{table}[htbp]
\centering
\small
\caption{CLIP scores reported by Zhuang et al.~\cite[Table~1]{zhuang2023}, as mean $\pm$ standard deviation.}
\label{tab:published}
\begin{tabular}{@{}lcc@{}}
\toprule
\textbf{Method} & \textbf{Untargeted} & \textbf{Targeted} \\
\midrule
No attack & $0.277\pm0.022$ & $0.229\pm0.030$ \\
Random & $0.271\pm0.021$ & $0.223\pm0.037$ \\
Greedy & $0.255\pm0.039$ & $0.204\pm0.037$ \\
Genetic & $0.203\pm0.042$ & $0.186\pm0.040$ \\
PGD & $0.226\pm0.041$ & $0.189\pm0.041$ \\
\bottomrule
\end{tabular}
\end{table}

\section{Method}

\subsection{Implementation}

The embedding is the final hidden state of the CLIP ViT-L/14 text encoder, flattened to 59,136 values (77 tokens $\times$ 768 dimensions) with padding tokens included, as in the authors' code~\cite{qfcode}. Suffixes are drawn from 62 characters, giving $62^{5}=916{,}132{,}832$ possibilities. Five search methods were run for each prompt and objective:
% AUTHOR CHECK: Verify the exact 62-character alphabet in the code.
% The recorded suffixes include punctuation and a middle dot; do not assume
% that "62 characters" means the usual alphanumeric alphabet.

\begin{itemize}
\item \textbf{Random (1):} a single random suffix.
\item \textbf{Greedy:} one sweep over the 5 positions from a random start, keeping the best of the 62 characters at each position.
\item \textbf{Genetic:} 10 rounds with 20 candidate suffixes per round.
The 2 best suffixes are kept unchanged. New candidates are created
by selecting the best of 3 randomly chosen candidates as each parent, combining parts of 2 parents at one split point, and applying random mutations with probability 0.3.
\item \textbf{Random (budget):} uniform random search over the same number of unique candidates as the genetic run for that prompt and objective.
\item \textbf{PGD:} the gradient-based attack from the paper, following the authors' code. Each of the 5 suffix slots holds a relaxed distribution over the 62 characters; Gumbel-softmax draws near-discrete samples, and a sign-gradient step on the relaxed logits maximises 1/cosine for 100 iterations. The final discrete suffix is scored with the same flattened-embedding cosine as the other methods, so its value is directly comparable. This ran on the GPU in float32 and made no diffusion calls. It was added as a reproduction of the paper's baseline, not as a stronger attack.
\end{itemize}

PGD operates on the text embedding only; no images were generated for its suffixes, so it appears in the text-embedding results (Table 3) but not in the image-score or detector tables.

\subsection{Scenes and image generation}

Two prompts from the paper's targeted examples were used. The first is ``A black bicycle against a brick wall'', with the bicycle as target and the wall as retained content. The second is ``A red apple on a plate'', with the plate as target and the apple as retained content.

Images were generated at $512\times512$ with Stable Diffusion v1.4 in half precision on a CUDA GPU on the University of Iowa Argon cluster, using 20 PNDM steps and a guidance scale of 7.5. Each scene used one seed (27), and the same initial noise was reused for all nine conditions: the clean prompt plus four methods under each of two objectives. This gave 18 images, and the safety checker, which was left enabled, flagged none of them.

Table~\ref{tab:setup} lists where this run differs from the paper.

\begin{table}[htbp]
\centering
\small
\caption{Differences between the paper's setup and this run.}
\label{tab:setup}
\begin{tabularx}{\linewidth}{@{}lYY@{}}
\toprule
\textbf{Setting} & \textbf{Paper~\cite{zhuang2023}} & \textbf{This run} \\
\midrule
Prompts & 20 & 2 \\
Images per prompt & 10 & 1 per condition \\
Genetic generations & 50 & 10 \\
Search methods & Random, greedy, genetic, PGD & Random, greedy, genetic, PGD, matched random \\
Mask sentences & Generated with ChatGPT & Ten fixed sentences defined in the code \\
\bottomrule
\end{tabularx}
\end{table}

\subsection{Evaluation}

Images were scored with the raw cosine similarity between CLIP image and text embeddings, without the rescaling sometimes applied to CLIP scores. Three texts were scored for each image: the original prompt, ``This is a photo of'' plus the target, and the same template plus the retained content. Results are reported as the drop from the clean image with the same seed, so a positive value means alignment fell and a negative value means it rose. Confidence intervals are omitted because there are only two scenes.

\subsection{Embedding analysis}

The 16 recorded suffixes were encoded again with the same encoder and revision, which reproduced every recorded cosine to within 0.001. Four measurements were then taken: the L2 distance from the clean embedding at each of the 77 token positions, the share of the total squared change that falls on padding positions, the cosine to a target-free prompt (``A brick wall'' or ``A red apple''), and the rank correlation across suffixes between text cosine and image score drop.

To give a reference distribution, 500 further random suffixes were encoded for each scene. These were used only for this analysis and are not part of any search budget. All values in the embedding figures are measured cosines or distances, so no dimensionality reduction is involved.

\section{Results}

\subsection{Text-embedding objective}

Table~\ref{tab:search} shows the search results. Lower cosine means a larger change in the embedding. The full and masked objectives use different sets of dimensions, so their values should not be compared directly.

\begin{table}[htbp]
\centering
\small
\caption{Best suffix and objective cosine for each scene, objective and method. Candidates is the number of unique suffixes evaluated.}
\label{tab:search}
\begin{adjustbox}{max width=\linewidth}
\begin{tabular}{@{}llllrr@{}}
\toprule
\textbf{Scene} & \textbf{Objective} & \textbf{Method} & \textbf{Suffix} & \textbf{Cosine} & \textbf{Candidates} \\
\midrule
Bicycle & Full & Random (1) & \texttt{\textasciicircum{}FF1\%} & 0.8628 & 1 \\
Bicycle & Full & Greedy & \texttt{FC\textperiodcentered{}!U} & 0.6722 & 306 \\
Bicycle & Full & Genetic & \texttt{\textbar{}ZPV\textgreater{}} & 0.6935 & 106 \\
Bicycle & Full & Random (budget) & \texttt{E2\textbar{},D} & 0.6901 & 106 \\
Bicycle & Full & PGD & \texttt{EE!\$7} & 0.7731 & 100 steps\\
\addlinespace
Bicycle & Masked & Random (1) & \texttt{.7P\textless{}Y} & 0.8510 & 1 \\
Bicycle & Masked & Greedy & \texttt{DJE\$8} & 0.6010 & 306 \\
Bicycle & Masked & Genetic & \texttt{61PV'} & 0.7411 & 105 \\
Bicycle & Masked & Random (budget) & \texttt{QMT\$9} & 0.7087 & 105 \\
Bicycle &Masked & PGD & \texttt{E4!\&U} & 0.6927 &100 steps\\
\addlinespace
Apple & Full & Random (1) & \texttt{H3X4\textbackslash{}} & 0.8254 & 1 \\
Apple & Full & Greedy & \texttt{.Y1\$)} & 0.7231 & 306 \\
Apple & Full & Genetic & \texttt{CT7\#S} & 0.6591 & 103 \\
Apple & Full & Random (budget) & \texttt{+\textbar{}SOX} & 0.7022 & 103 \\
Apple & Full & PGD & \texttt{G63S\textbar{}} &0.5154 &100 steps\\
\addlinespace
Apple & Masked & Random (1) & \texttt{\#5)CT} & 0.7577 & 1 \\
Apple & Masked & Greedy & \texttt{GP38I} & 0.5953 & 306 \\
Apple & Masked & Genetic & \texttt{Y!KD8} & 0.7107 & 94 \\
Apple & Masked & Random (budget) & \texttt{X27QH} & 0.7457 & 94 \\
Apple & Masked & PGD & \texttt{Q7\textbar{}MG} & 0.5389 & 100 steps \\
\bottomrule
\end{tabular}
\end{adjustbox}
\end{table}

Every search beat a single random suffix. Genetic search beat matched random search on both objectives for the apple scene, and lost on both for the bicycle scene. Among the discrete searches, greedy reached the lowest masked cosine on both scenes, but it evaluated about three times as many candidates as the genetic runs, so this is not a like-for-like comparison.

PGD reached the lowest cosine of any method on both apple objectives (0.5154 full and 0.5389 masked). On the bicycle scene it placed second on the masked objective, behind greedy (0.6927 against 0.6010), and fourth on the full objective (0.7731), ahead of only the single random suffix. Its outcome also depends on the random Gumbel samples. An earlier run of the same code on a CPU, with the same seed but a different random-number stream, gave 0.5886, 0.6205, 0.6064 and 0.8303 for the four cases in table order, and PGD's rank among the five methods changed in two of the four cases. Both runs agree with the paper's finding that PGD is comparable to the discrete searches rather than uniformly stronger. Its cosines are not directly comparable to the paper's Table 1 numbers, which are image CLIP scores rather than text-embedding cosines.

The masks selected 23,310 of 59,136 dimensions (39.4\%) for ``a black bicycle'' and 19,373 (32.8\%) for ``a plate''. The search with the four discrete methods used 2,117 text encodings, made no diffusion calls and took 5.3 seconds on the CUDA GPU used for this run.
% AUTHOR CHECK: Table~\ref{tab:search} totals 2,044 unique candidates.
% Clarify what additional encodings are included in the stated total of 2,117.

\subsection{Image scores}

Table~\ref{tab:image-scores} gives the mean drops. The untargeted rows are scored against the original prompt and the targeted rows against the target phrase.

\begin{table}[htbp]
\centering
\small
\caption{Mean paired drop in CLIP score over the two scenes. Retained drop is the same measure for the retained content.}
\label{tab:image-scores}
\begin{tabular}{@{}llrr@{}}
\toprule
\textbf{Objective} & \textbf{Method} & \textbf{Mean drop} & \textbf{Retained drop} \\
\midrule
Untargeted & Random (1) & $+0.0002$ & $+0.0091$ \\
Untargeted & Greedy & $-0.0240$ & $-0.0078$ \\
Untargeted & Genetic & $-0.0345$ & $+0.0059$ \\
Untargeted & Random (budget) & $+0.0438$ & $+0.0223$ \\
\addlinespace
Targeted & Random (1) & $-0.0151$ & $-0.0061$ \\
Targeted & Greedy & $-0.0041$ & $+0.0094$ \\
Targeted & Genetic & $+0.0093$ & $+0.0518$ \\
Targeted & Random (budget) & $-0.0170$ & $-0.0098$ \\
\bottomrule
\end{tabular}
\end{table}

Untargeted genetic search raised alignment with the original prompt on average, and the largest untargeted drop came from matched random search. Targeted genetic search was the only targeted method with a positive target drop, but its retained content dropped by more. Most of that retained drop came from the bicycle scene, where the wall score fell by 0.1045 while the target score fell by only 0.0113. The lowest text cosine also did not produce the largest image effect: among the methods with generated images, targeted greedy had the lowest masked cosine on both scenes, yet its mean target drop was slightly negative.

\subsection{Visual inspection}

Figures~\ref{fig:bicycle} and~\ref{fig:apple} show the generated images. The top row uses the full objective and the bottom row the masked objective, and all images in a figure share the same initial noise. These are informal, unblinded observations and no labels were collected.

The suffixes clearly change composition and style. Several bicycle images shift from a photograph towards an illustration, and in the apple scene the fruit changes in number and size. The untargeted genetic image replaces the real fruit with an apple motif on tableware. Removal did not occur in the targeted images, however. The bicycle remains visible in all four targeted bicycle images, and a plate remains visible in all four targeted apple images. Section~\ref{sec:detector} checks this with an object detector instead of by eye.

\begin{figure}[!h]
\centering
\reportfigure{fig_images_bicycle.png}
\caption{Generated images for ``A black bicycle against a brick wall'', seed 27.}
\label{fig:bicycle}
\end{figure}

\begin{figure}[!h]
\centering
\reportfigure{fig_images_apple.png}
\caption{Generated images for ``A red apple on a plate'', seed 27.}
\label{fig:apple}
\end{figure}

\subsection{Embedding analysis}

\textbf{Where the change happens:} Figure~\ref{fig:token-change} follows the greedy masked suffix through the encoder, one token at a time. Each position in CLIP's text encoder can attend only to itself and earlier tokens, so a suffix cannot affect anything before it. The measurements confirm this up to numerical precision: across all 16 suffixes, the largest change at any position before the suffix is 0.00003, against distances of about 12 to 37 from the suffix onwards. The tokens ``bicycle'' and ``plate'' therefore reach the diffusion model effectively unchanged.

Most of the change falls on padding, which fills the remaining 63 to 66 of the 77 positions. Across the 16 suffixes, padding holds 72.2\% to 92.7\% of the squared change. It also holds about 87\% of the squared norm of the clean embedding (86.8\% and 87.9\%), so the objective largely measures change in padding.

\textbf{What the mask selects:} The lower row of Figure~\ref{fig:token-change} shows the same pattern in the mask. Of the 23,310 dimensions selected for ``a black bicycle'', 99.14\% are on padding positions, and of the 19,373 selected for ``a plate'', 99.51\% are. No dimension is selected at positions 0 to 6, because the vote must be unanimous and the embedding there does not change for every context sentence when the target phrase is added. The mask therefore contains nothing at the target word's own position in either prompt.

\begin{figure}[!h]
\centering
\reportfigure{fig_token_change.png}
\caption{One suffix per scene, token by token. The upper row shows how far each position moves from the clean embedding, and the lower row shows how much of that position the mask selects. The mask is fixed by position and does not depend on the suffix. The padding bar is the mean over the remaining positions.}
\label{fig:token-change}
\end{figure}

\textbf{Whether the attack moves towards a target-free prompt:} A suffix that removed the target might be expected to move the embedding closer to the same prompt without the target. Figure~\ref{fig:similarity} plots both similarities directly, and Table~\ref{tab:targeted} gives the targeted values. Every suffix moves the other way. On the bicycle scene, the clean prompt has cosine 0.6360 to ``A brick wall'', and the eight attacked prompts range from 0.3988 to 0.5414. On the apple scene, the clean prompt has cosine 0.7470 to ``A red apple'', and the attacked prompts range from 0.5168 to 0.6284. None of the 500 random suffixes per scene came closer than the clean prompt either. This suggests that the drop is not specific to the search methods.

\begin{figure}[!h]
\centering
\reportfigure{fig_similarity_plane.png}
\caption{Each point is one prompt, placed by its measured cosine to the clean prompt and to the target-free prompt. The horizontal line marks the clean prompt's value, and points above it would be closer to the target-free prompt. The labelled points are the suffixes shown in Figure~\ref{fig:token-change}.}
\label{fig:similarity}
\end{figure}

% Two panels keep the original ten-column table readable.
\begin{table}[!h]
\centering
\small
\setlength{\tabcolsep}{4pt}
\caption{Targeted-attack examples with measured values. Panel (a) reports embedding similarities; \emph{Tokens} gives the number of suffix tokens. Cosine to target-free compares the full embedding with ``A brick wall'' or ``A red apple''. Panel (b) reports the share of squared change on padding positions and the paired image-score drops.}
\label{tab:targeted}
\textbf{(a) Suffixes and embedding similarities}\par\smallskip
\begin{adjustbox}{max width=\linewidth}
\begin{tabular}{@{}lllrrrr@{}}
\toprule
\textbf{Scene} & \textbf{Method} & \textbf{Suffix} & \textbf{Tokens} & \shortstack{\textbf{Masked}\\\textbf{cos.}} & \shortstack{\textbf{Full}\\\textbf{cos.}} & \shortstack{\textbf{Cos. to}\\\textbf{target-free}} \\
\midrule
Bicycle & Clean & none & 0 & 1.0000 & 1.0000 & 0.6360 \\
Bicycle & Random (1) & \texttt{.7P\textless{}Y} & 5 & 0.8510 & 0.8608 & 0.5009 \\
Bicycle & Greedy & \texttt{DJE\$8} & 4 & 0.6010 & 0.6615 & 0.3988 \\
Bicycle & Genetic & \texttt{61PV'} & 4 & 0.7411 & 0.7684 & 0.4805 \\
Bicycle & Random (budget) & \texttt{QMT\$9} & 4 & 0.7087 & 0.7550 & 0.4479 \\
\addlinespace
Apple & Clean & none & 0 & 1.0000 & 1.0000 & 0.7470 \\
Apple & Random (1) & \texttt{\#5)CT} & 4 & 0.7577 & 0.7642 & 0.6218 \\
Apple & Greedy & \texttt{GP38I} & 4 & 0.5953 & 0.6095 & 0.5168 \\
Apple & Genetic & \texttt{Y!KD8} & 4 & 0.7107 & 0.7160 & 0.5421 \\
Apple & Random (budget) & \texttt{X27QH} & 5 & 0.7457 & 0.7395 & 0.5572 \\
\bottomrule
\end{tabular}
\end{adjustbox}

\medskip
\textbf{(b) Padding contribution and image-score drops}\par\smallskip
\begin{tabular}{@{}llrrr@{}}
\toprule
\textbf{Scene} & \textbf{Method} & \textbf{Padding share} & \textbf{Target drop} & \textbf{Retained drop} \\
\midrule
Bicycle & Clean & n/a & n/a & n/a \\
Bicycle & Random (1) & 72.2\% & $-0.0060$ & $+0.0107$ \\
Bicycle & Greedy & 87.7\% & $+0.0141$ & $+0.0377$ \\
Bicycle & Genetic & 85.8\% & $+0.0113$ & $+0.1045$ \\
Bicycle & Random (budget) & 87.1\% & $-0.0079$ & $+0.0167$ \\
\addlinespace
Apple & Clean & n/a & n/a & n/a \\
Apple & Random (1) & 87.9\% & $-0.0241$ & $-0.0229$ \\
Apple & Greedy & 89.1\% & $-0.0223$ & $-0.0189$ \\
Apple & Genetic & 87.4\% & $+0.0073$ & $-0.0008$ \\
Apple & Random (budget) & 82.3\% & $-0.0261$ & $-0.0363$ \\
\bottomrule
\end{tabular}
\end{table}

\textbf{How the searches compare with random suffixes:} Figure~\ref{fig:random} compares each way of choosing a suffix. The mean full cosine for a random suffix is 0.8647 on the bicycle scene and 0.8493 on the apple scene, and the searches all reach well below the typical range. Random sampling also gets there when given more attempts, since on the apple scene the best of 500 suffixes is lower than every suffix from the discrete searches, though not lower than PGD's (0.5154 and 0.5389). That sample is larger than the budgets of the multi-candidate search methods (94 to 306).

\begin{figure}[!h]
\centering
\reportfigure{fig_searches_vs_random.png}
\caption{Objective cosine for each way of choosing a suffix, with the number of suffixes tried. Compare each point with the band of the same colour, since the full and masked objectives use different dimensions.}
\label{fig:random}
\end{figure}

\textbf{Whether text change predicts image change:} Figure~\ref{fig:text-image} plots the text cosine against the image score drop for all 16 suffixes. If a larger text change produced a larger image effect, the points would slope down from upper left to lower right, giving a clearly negative rank correlation. The values are $+0.07$ for the full objective and $-0.29$ for the masked objective. Only the masked value points in the expected direction, and it is weak. With eight suffixes in each panel, this is only descriptive.

\begin{figure}[!h]
\centering
\reportfigure{fig_text_vs_image.png}
\caption{Text cosine against image score drop for each suffix, using the individual image-score drops whose means are reported in Table~\ref{tab:image-scores}.}
\label{fig:text-image}
\end{figure}

\subsection{Separate object-detection check}
\label{sec:detector}

The CLIP scores rely on the same model family used in the attack. As a separate check on those scores and the visual inspection, the 18 images were scored with OWL-ViT, an open-vocabulary object detector~\cite{minderer2022}, which was not used anywhere in the attack. For each image, the detector reported confidence scores for the target object and the retained object. A detection threshold of 0.10 was fixed in advance.

The detector found the target object in every image. In the two clean images, its confidence is 0.686 for the bicycle and 0.361 for the plate. In all eight targeted attacks and all eight untargeted attacks, the target stayed above the threshold. Thus, no attack met this detector-based criterion for object removal. Figure~\ref{fig:detector} shows the confidences, and Table~\ref{tab:detector} gives the mean change from the clean image.

\begin{figure}[!h]
\centering
\reportfigure{fig_detector.png}
\caption{Detector confidence for the target object (circles) and the retained object (crosses) in each image. Points to the right of the dashed line are counted as detected. Every target circle is to the right of it.}
\label{fig:detector}
\end{figure}

\begin{table}[!h]
\centering
\small
\caption{Change in the detector's target confidence relative to the clean image, averaged over the two scenes. A positive drop means the target was found less confidently. Targets removed counts images where the target was above the threshold in the clean image but below it after the attack.}
\label{tab:detector}
\begin{tabular}{@{}llrc@{}}
\toprule
\textbf{Objective} & \textbf{Method} & \shortstack{\textbf{Mean target}\\\textbf{confidence drop}} & \shortstack{\textbf{Targets}\\\textbf{removed}} \\
\midrule
Untargeted & Random (1) & $+0.1098$ & 0 of 2 \\
Untargeted & Greedy & $+0.0125$ & 0 of 2 \\
Untargeted & Genetic & $+0.0630$ & 0 of 2 \\
Untargeted & Random (budget) & $+0.2339$ & 0 of 2 \\
\addlinespace
Targeted & Random (1) & $+0.0465$ & 0 of 2 \\
Targeted & Greedy & $-0.0049$ & 0 of 2 \\
Targeted & Genetic & $+0.1297$ & 0 of 2 \\
Targeted & Random (budget) & $+0.0172$ & 0 of 2 \\
\bottomrule
\end{tabular}
\end{table}

The mean confidence drops range from $-0.0049$ to 0.2339. For some individual images, the detector was more confident after the attack. Confidence scores themselves do not depend on the chosen threshold; the binary detection and removal decisions do. At the fixed 0.10 threshold, the target was detected after every attack in this run. OWL-ViT uses a CLIP-based backbone, so this provides a separate evaluation pipeline but not a fully independent model family. The retained-object confidence is low even in the clean apple-scene image (0.137), so the crosses in Figure~\ref{fig:detector} should be read only as a rough guide.

\section{Discussion}

\subsection{Comparison with the paper}

The first claim is partly supported. Short suffixes substantially reduced the text-space objective with no diffusion queries, and they visibly changed the images. The CLIP score drops were much smaller and less consistent than the 0.074 reported for the paper's genetic attack, although the two figures come from very different sample sizes.

The second claim was not reproduced. The paper shows the bicycle and plate being removed from these same prompts~\cite[Fig.~4]{zhuang2023}, whereas neither was removed in this run. This run is too limited to contradict the paper, since it used fewer generations, fewer sampling steps, one seed and different mask sentences. The separate detector in Section~\ref{sec:detector} agrees with the visual inspection that neither object was removed at the chosen detection threshold. This check does not use the report's CLIP similarity score, although the detector shares the CLIP model family. The results show that minimizing the masked objective did not guarantee targeted removal in these cases.

The matched control suggests that the comparison between optimizers needs care. On these prompts, random search with the same budget was competitive with a 10-generation genetic algorithm. Applying the same control across the full 20-prompt setting would show whether the paper's ranking comes from the search strategy or the number of evaluations.

The embedding analysis suggests an explanation for the failed removal, which this run does not test directly. The target token's embedding is effectively unchanged before and after the attack, so the diffusion model's cross-attention still receives it. The mask selects almost only padding positions, and the search moves the embedding away from the target-free prompt as well as from the clean prompt. Testing this would require measuring image outcomes after directly changing the target token's embedding or the padding embeddings. The authors' code flattens the embedding in the same way, so the same concentration on padding may apply to the original experiments, but this was only measured on the two prompts here.

\subsection{Limitations}

These results have several limitations.

\begin{itemize}
\item Two scenes and one seed are too few for statistical conclusions, and each method had one search run per case.
\item The 20-step sampling and 10-generation search are shorter than the paper's settings, which limits direct comparison.
\item The attack objective and the image evaluation both rely on CLIP, so the evaluator may share sensitivities with the attacked encoder. The evaluation texts do not contain the suffix. The detector in Section~\ref{sec:detector} provides a separate check, but it still has a CLIP-based backbone, and its binary detection decisions depend on the chosen threshold.
\item The masks depend on the chosen sentences, and with about a third of all dimensions selected they were not especially selective.
\item The cosine to a target-free prompt uses one reference prompt per scene and is itself dominated by padding positions.
\item The measured run used one CUDA GPU with half-precision weights. An earlier run on Apple MPS found the same discrete-search suffixes, but its images differed slightly, its mean image score drops differed by up to 0.014, and its PGD suffixes were different. The notebook has not been executed on hosted Google Colab.
\end{itemize}

Blinded human labels or a detector with no shared backbone would strengthen the evaluation further, especially for the retained content.

\section{Conclusion}

This reproduction confirms that five-character suffixes can lower the CLIP text-embedding objective substantially without querying the diffusion model. Those changes altered the images but did not reduce CLIP alignment reliably. A separate detector found the target object in every image at the chosen threshold, and the genetic algorithm did not consistently outperform random search with the same budget. The embedding analysis shows that the change falls mainly on padding positions and leaves the target token effectively unchanged, which may help explain why objects were not removed. A useful next step is to run the implementation's extended profile (20 scenes, 50 generations, 50 steps and ten seeds) and score it with a detector that does not share the CLIP backbone.

\begin{thebibliography}{9}
\bibitem{zhuang2023}
H. Zhuang, Y. Zhang, and S. Liu.
A pilot study of query-free adversarial attack against Stable Diffusion.
In \emph{CVPR Workshops}, 2023. arXiv:2303.16378.

\bibitem{rombach2022}
R. Rombach, A. Blattmann, D. Lorenz, P. Esser, and B. Ommer.
High-resolution image synthesis with latent diffusion models.
In \emph{CVPR}, 2022.

\bibitem{radford2021}
A. Radford, J. W. Kim, C. Hallacy, et al.
Learning transferable visual models from natural language supervision.
In \emph{ICML}, 2021.

\bibitem{qfcode}
OPTML-Group. QF-Attack.
\url{https://github.com/OPTML-Group/QF-Attack}

\bibitem{minderer2022}
M. Minderer, A. Gritsenko, A. Stone, et al.
Simple open-vocabulary object detection with vision transformers.
In \emph{ECCV}, 2022.
\end{thebibliography}

\section*{Appendix: Software and models}
\phantomsection
\label{sec:software}

Python 3.10.8, torch 2.6.0, transformers 4.49.0, diffusers 0.32.2, accelerate 1.4.0, huggingface-hub 0.28.1 and numpy 2.0.2, on a CUDA GPU on the University of Iowa Argon cluster (Linux). Search seed: 2026.

Object detection used OWL-ViT (\nolinkurl{google/owlvit-base-patch32}).

Stable Diffusion v1.4 revision:\\
\nolinkurl{133a221b8aa7292a167afc5127cb63fb5005638b} (fp16 weights).

CLIP ViT-L/14 revision:\\
\nolinkurl{32bd64288804d66eefd0ccbe215aa642df71cc41}.

The code, measured results and full-resolution images accompany this report.

\end{document}